\documentclass[12pt]{article}

% --- Packages (The "plugins" of LaTeX) ---
\usepackage[utf8]{inputenc} % Character encoding
\usepackage{geometry}       % Page margins
 \geometry{margin=1in}
\usepackage{graphicx}       % For inserting images
\usepackage{amsmath}        % For math equations
\usepackage{listings}
\usepackage{xcolor} % Required for syntax highlighting colors

% Define colors for Python syntax
\definecolor{codegreen}{rgb}{0,0.6,0}
\definecolor{codegray}{rgb}{0.5,0.5,0.5}
\definecolor{codepurple}{rgb}{0.58,0,0.82}
\definecolor{backcolour}{rgb}{0.95,0.95,0.92}

\lstdefinestyle{mystyle}{
    backgroundcolor=\color{backcolour},   
    commentstyle=\color{codegreen},
    keywordstyle=\color{magenta},
    numberstyle=\tiny\color{codegray},
    stringstyle=\color{codepurple},
    basicstyle=\ttfamily\footnotesize,
    breakatwhitespace=false,         
    breaklines=true,                 
    captionpos=b,                    
    keepspaces=true,                 
    numbers=left,                    
    numbersep=5pt,                  
    showspaces=false,                
    showstringspaces=false,
    showtabs=false,                  
    tabsize=2
}

\lstset{style=mystyle}

% --- Document Information ---
\title{DE 414: Practical 1 Report}
\author{Amy McDermott (26911264)}
\date{\today}

\begin{document}

\maketitle

\section*{Introduction}
This report focuses on documenting fundamental techniques in the Python programming language. This includes File I/O, NumPy array creation, indexing, slicing and visualization with Matplotlib.

\section{Question 1}
Firstly, the essential Python libraries \texttt{numpy} and \texttt{matplotlib} are imported.

\begin{lstlisting}[language=Python, caption={Import libraries}]
import numpy as np
import matplotlib.pyplot as plt
\end{lstlisting}

\section{Question 2}
The MNIST dataset is read from their respective files and stored as \texttt{numpy} variables. 

\begin{lstlisting}[language=Python, caption={Read the MNIST dataset into numpy variables}]
# Define an empty list to store the training data
data = []

# Read the data line by line into the list
with open("data/train.csv") as f:
    for line in f:
        data.append(line.rstrip().split(","))

# Convert the list into a numpy array
data = np.array(data, dtype = np.int32)

# Slice data to seperate training from target data
train = data[:,1:]
train_target = data[:,0]

# Repeat this process for the test data
test_data = []

with open("data/test.csv") as f:
    for line in f:
        test_data.append(line.rstrip().split(","))

test_data = np.array(test_data, dtype = np.int32)

# Slice the test data into raw data and its corresponding targets
test = test_data[:,1:785]
test_target = test_data[:,0]
\end{lstlisting}

\section{Question 3}
As defined in the practical 1 documentation, the variables must result with the following shapes:
\begin{enumerate}
    \item Train: $(60000\times784)$
    \item Train targets:  $(60000\times1)$
    \item Test: $(10000\times784)$
    \item Test targets: $(10000\times1)$
\end{enumerate}
This is verified using the following code:

\begin{lstlisting}[language=Python, caption={Confirm the shapes of the training and test data}]
print(train.shape) # Prints the dimensions of the training data to the terminal
print(train_target.shape)
print(test.shape)
print(test_target.shape)
\end{lstlisting}
The terminal output from the above listing, confirms that the training and test data was correctly read into \texttt{numpy} variables.

\section{Question 4}
The data is then normalised by means of linear scaling so that it lies within the range [0, 1]. The MNIST pixels are integers in the range 
[0, 255]. Dividing by 255.0 scales the values to the required range.

\begin{lstlisting}[language=Python, caption={Normalise the data using linear scaling}]
train_norm = train / 255.0 # Diving by a float prevents integer division
test_norm = test / 255.0

# Quick sanity checks to confirm the normalization worked as expected
min_train = np.min(train_norm)
max_train = np.max(train_norm)
min_test = np.min(test_norm)
max_test= np.max(test_norm)

print(f"The min of the training data is {min_train} and the maximum is {max_train}")
print(f"The min of the test data is {min_test} and the maximum is {max_test}")
\end{lstlisting}
The output from the above listing confirms that the minimum and maximum values from both the training and test data lie within the [0,1] interval.

\section{Question 5}
Plot the first example from the training set using \texttt{matplotlib}.

\begin{lstlisting}[language=Python, caption={Plot an example from the training data}]
train_ex = train_norm[0,:] # Select the data from the first row
train_ex = np.reshape(train_ex, (28, 28)) # Reshape the pixel data so that it can be correctly plotted

plt.imshow(train_ex) # Plot the data
plt.show()
\end{lstlisting}

\section{Question 6}
Weights from a logistic regression model are loaded and stored in a \texttt{numpy} array.

\begin{lstlisting}[language=Python, caption={Load weights into a numpy array}]
weights = [] # Create an empty list

# Read in the weights line by line into the list
with open("data/weights.csv", 'r') as f:
    for line in f:
        weights.append(line.rstrip().split(","))

# Convert the weights into a numpy array
weights = np.array(weights, dtype = np.float32)
\end{lstlisting}

\section{Question 7}
Using the \texttt{SoftmaxRegression} model, a prediction is made for an instance of data from the test set at index 7.

\begin{lstlisting}[language=Python, caption={Predict digit of example from the test set}]
from models import SoftmaxRegression
model = SoftmaxRegression()

model.load(weights)

# Extract an example from test set
example = test_norm[7,:]
example = example.reshape((784,1)).T

# Returns a (1,10) array of probabilities. This row vector of length 10 contains the predicted probabilities for each digit 
prob = model(example)
max_index = np.argmax(prob) #To obtain the predicted class id use `np.argmax(prob)` which returns the index of the maximum probability.

print(f"The predicted value of the mnist digit example is {max_index}") # Print out the predicted value.
\end{lstlisting}

\section{Question 8}
A bar graph is plotted to display the model probabilities computed for each possible output digit for the test set sample from index 7.

\begin{lstlisting}[language=Python, caption={Display a bar graph of probabilities of each digit for a test set example}]
x_axis = np.arange(0, 10) # Create array of evenly spaced values between 0 and 10
plt.bar(range(10), prob[0]) # Pass in x-axis positions and corresponding values
plt.show()
\end{lstlisting}

\section{Question 9}
A Python function is defined to calculate the model accuracy. This function is used to calculate the average accuracy on the test set.

\begin{lstlisting}[language=Python, caption={Compute accuracy of model classifier}]
# Function to determine accuracy
def accuracy(num_correct, total_pred):
    return num_correct/total_pred # Divide the number of total correct classifications by total predictions

test_pred = [] # Define an empty list to store the test data classification results

# Loop through every data instance (row)
for row in test_norm:
    row = row.reshape((784,1)).T
    probs = model(row) # Returns an array with 10 indices, storing the probability of the instance being classified as digit i
    max_index = np.argmax(probs) # Classification result determined using index of array element containing the maximum probability

    # Append new row of probabilities for next digit
    test_pred.append(max_index)

test_pred = np.array(test_pred) # Convert the list to a numpy array

eq = np.equal(test_pred, test_target) # Test whether our model classifications are equal to their corresponding target values
total_correct = np.sum(eq) # Sum the total correct classifications

acc = accuracy(total_correct, test_target.size) # Use the function to determine accuracy of classifier
print(acc) # Print accuracy to the terminal
\end{lstlisting}
This results in a model accuracy of 0.9229. Thus, 92.3\% of all instances of data that are classified using this model, are correctly identified as their corresponding digit.

\end{document}